\begin{papercard}{pangakis2023}
\cf{Stages} \stagetag{L} (workflow for LLM annotation with human validation), \stagetag{C} (codebook as prompt; human-in-the-loop codebook refinement), \stagetag{I} (held-out comparison with expert labels; consistency scores). A 2023 preprint by three political-science and communication researchers that argues ``any automated annotation process using an LLM must validate the LLM's performance against labels generated by humans'' and tests a five-step workflow by replicating 27 annotation tasks from 11 social-science datasets with GPT-4 \pc{pangakis2023}{1}{we argue that any automated annotation process using an LLM must validate the LLM’s performance against labels generated by humans}. It is the empirical anchor for ``validate task by task'' that Törnberg (2024) and Baumann et al.\ (2025) build on.

\cf{Motivation} Human annotators (crowd-workers or undergraduates) suffer from limited attention, fatigue and changing perceptions of categories during labelling, which can cause errors that are ``unobservable and correlated''; LLMs are faster, cheaper and reproducible \pc{pangakis2023}{2}{often suffer from a limited attention span, fatigue, and changing perceptions of the underlying conceptual categories}. Prior evidence is mixed (Gilardi et al.\ positive; Reiss negative). Existing studies report minimal metrics (only accuracy or intercoder agreement), cover few insular tasks and use public benchmark datasets that may be contaminated by training data \pc{pangakis2023}{3}{these tests are plausibly affected by contamination}. Main argument: researchers must always validate against a subset of \emph{high-quality, subject-expert} human labels \pc{pangakis2023}{3}{a researcher using an LLM for automated annotation must always validate the LLM’s performance against a subset of high-quality human-annotated labels}, because failures arise from prompt quality, noisy texts and conceptual difficulty and will persist as models improve.

\cf{The five-step workflow (Figure 1)} (1) Researcher writes task-specific instructions (a codebook). (2) Subject-matter experts annotate a random subset of texts using it. (3) The LLM annotates a subset of the human-labelled data with the same codebook and is evaluated against the human labels. (4) If performance is low, refine the codebook to emphasise the misclassifications, repeating steps 2--3 if needed. (5) Using the final codebook, test the LLM on the \emph{remaining} human-labelled samples; performance on this held-out data decides whether the LLM can be used \pc{pangakis2023}{6}{Performance on this held-out data should be used to determine whether an LLM can be effectively utilized for automated annotation}. Requirements: at least two subject experts and the LLM annotate the same units; \emph{the same codebook} serves as human instructions and LLM prompt \pc{pangakis2023}{6}{Critically, both the human coders and the LLM should use the same codebook for annotation, where the codebook serves as the LLM’s prompt instructions}; metrics accuracy, precision, recall, F1; suggested sample 250--1,250 random texts, more for rare cases below 1\,\% prevalence \pc{pangakis2023}{6}{the number of samples annotated should range between 250 and 1,250 random text samples}. Re-labelling the tuning subset after the codebook update measures how much the edit changed model behaviour but ``is overfitting to a particular subset of data and thus should not be taken as the barometer of performance'' \pc{pangakis2023}{6}{This exercise is overfitting to a particular subset of data and thus should not be taken as the barometer of performance}.

\cf{Replication design} 27 tasks from 11 datasets in recent articles from journals such as APSR, AJPS, PNAS, Science Advances, JPSP; each original article's human-coded data serve as ground truth; to avoid contamination the datasets come from password-protected archives (e.g.\ Dataverse) or directly from authors \pc{pangakis2023}{7}{To avoid the potential for contamination, we rely exclusively on datasets stored in password-protected data archives}; the original codebook when available, otherwise quotes or paraphrases from the article. Every dimension is a separate binary task \pc{pangakis2023}{7}{we treat every dimension as a separate binary annotation task}. Over 200,000 texts classified with the GPT-4 API at a cost of about \$420; roughly 2--3 hours for seven iterations on 1,000 texts \pc{pangakis2023}{7}{Across all 27 tasks, we annotate slightly over 200,000 text samples using OpenAI’s GPT-4 API}. Codebook updates: at most one round per task. Results use held-out samples only.

\cf{Results} Across the 27 tasks: accuracy median .850 (range .674--.981), precision median .650 (.033--.957), recall median .829, F1 median .707 (.059--.969) \pc{pangakis2023}{4}{Across all tasks, we report a median accuracy of 0.850 and a median F1 of 0.707}. Recall exceeds precision in 20 of 27 tasks \pc{pangakis2023}{8}{LLM classification performance is stronger in recall than precision for 20 of the 27 tasks}; eight tasks have both above .7; nine of 27 have precision or recall below .5 and three have both below .5, so ``for a full one-third of tasks'' the LLM misses half the positives or produces more false than true positives \pc{pangakis2023}{8}{for a full one-third of tasks, the LLM either missed at least half of the true positive cases, had more false positives than true positives, or both}. Performance varies even within one dataset (F1 .259 vs.\ .811 for two tasks of Card et al.\ 2022) \pc{pangakis2023}{10}{In the most extreme example, F1 ranged from 0.259 to 0.811 on two separate tasks within Card et al. (2022), a difference of 0.552}. Whether the original codebook was available made no observable difference (footnote). Codebook edits gave modest gains in accuracy and F1 (precision up in most cases, recall down in more cases than up), so ``the subtleties of prompt construction do not appear to be a significant lever on performance'' \pc{pangakis2023}{12}{the subtleties of prompt construction do not appear to be a significant lever on performance}, though refinement can rule out unclear instructions.

\cf{Consistency score} Classify each text repeatedly at temperature above 0 (the authors recommend at least three runs; they use temperature 0.6 and seven runs); the score is the share of runs that match the modal label. Accuracy is 19.4 percentage points higher for texts with a score of 1.0 than for those below 1.0; true-positive rate +16.4 and true-negative rate +21.4 points; 85.1\,\% of texts are fully consistent, so a score below 1 flags edge cases \pc{pangakis2023}{10}{accuracy is 19.4 percentage points higher for text samples labeled with a consistency score of 1.0 compared with those labeled with a consistency score less than 1.0}. The 0.6 value is ``arbitrary'' but repeated sampling beat a single temperature-0 run \pc{pangakis2023}{10}{While our choice of 0.6 is arbitrary}.

\cf{Four use cases (Table 2)} (1) check the quality of human-labelled data (low LLM--human agreement means that humans, LLM or both erred); (2) prioritise cases for human review (manually review items with consistency below 1.0; with high recall, let the LLM propose positives for human confirmation) \pc{pangakis2023}{14}{A researcher can manually review annotations with consistency scores lower than 1.0}; (3) produce labelled training data for a supervised classifier; (4) classify the whole corpus directly. If performance is poor, refine, abandon LLM use for that dimension, or disaggregate a complex dimension into component indicators \pc{pangakis2023}{15}{they may achieve better performance by disaggregating “hate speech” into component indicators}.

\cf{Conclusion} Heterogeneity is large across tasks and even within a dataset; ``it is imperative that researchers validate the performance of LLMs on a task-by-task basis'' \pc{pangakis2023}{15}{it is imperative that researchers validate the performance of LLMs on a task-by-task basis, as we find significant heterogeneity in performance, even across tasks within a single dataset}. Code released as a Python package (gpt\_annotate).

\cf{Limits (card)} Preprint of June 2023 with an unspecified GPT-4 snapshot that is no longer available; the ground truth is the original study's human labels, whose reliability is not re-checked and is not reported per task (so human--human agreement is not the benchmark); binary decomposition means class imbalance drives precision; F1 and accuracy but no chance-corrected $\kappa$ or $\alpha$ and no confidence intervals; sample-size range 250--1,250 given without derivation; at most one codebook-update round; no analysis of why tasks fail; no costs for expert human labelling; all English social-science text.

\cf{Use for this thesis} (1) Gives the thesis a citable five-step, human-first protocol for stage L that maps to the thesis stages: codebook (C) $\to$ two experts code a random subset (I) $\to$ LLM on tuning subset $\to$ refine $\to$ untouched held-out test. (2) The same-codebook rule: the LLM prompt is the codebook text, so the codebook must be written for both readers. (3) Consistency scoring is a cheap way to flag units for human review; it is suited to the thesis where an LLM is only a screening aid. (4) The finding that one-third of tasks fail even with GPT-4 is a reason to pre-register the minimum acceptable agreement with the human consensus before looking at held-out results, and to report per-code results. (5) Pair with Törnberg for the procedure and Baumann for the downstream-inference risk.
\end{papercard}
